% Generated draft: requires mathematical and source review.
\documentclass[11pt]{article}
\usepackage[margin=25mm]{geometry}
\usepackage{fontspec}
\setmainfont{Latin Modern Roman}
\usepackage{amsmath,amssymb,mathtools,unicode-math}
\setmathfont{Latin Modern Math}
\usepackage{enumitem,needspace,xurl,hyperref}
\setlength{\parindent}{0pt}
\setlength{\parskip}{5pt}
\setlength{\emergencystretch}{3em}
\newcommand{\N}{\mathbb{N}}
\newcommand{\Z}{\mathbb{Z}}
\newcommand{\Q}{\mathbb{Q}}
\newcommand{\R}{\mathbb{R}}
\newcommand{\C}{\mathbb{C}}
\newcommand{\F}{\mathbb{F}}
\newcommand{\A}{\mathbb{A}}
\newcommand{\PP}{\mathbb{P}}
\newcommand{\E}{\mathbb{E}}
\newcommand{\eps}{\varepsilon}
\newcommand{\dd}{\,\mathrm{d}}
\begin{document}
\section*{UnsolvedMath ID: 10600023}
\textbf{Source problem code:} catalogue entry
\subsection{Definitions and mathematical statement}
For a link \(L\), let \(V_L(t)\) denote the Jones polynomial and \(\operatorname{Kh}^{i,j}(L;R)\) its Khovanov homology, with homological grading \(i\), quantum grading \(j\), and coefficient system \(R\). In the usual classical setting, the Jones polynomial is a skein invariant and Khovanov homology is obtained from a cube-of-resolutions complex whose appropriately normalized graded Euler characteristic recovers the Jones polynomial. Framed/oriented, reduced/unreduced and grading conventions must be fixed before stating an exact formula; the record gives none.

Problem-list item 23 asks for \emph{new geometric or topological interpretations} of the Jones polynomial and of Khovanov homology. Both named objects are targets of the request. The source does not select a link class, a geometric construction, an identity, a realization category, or a notion of equivalence by which an interpretation is to be judged. In particular the context of a virtual-knot survey does not by itself specify that a particular virtual extension of either invariant is intended.

The word ``new'' requires comparison with an unstated body of existing interpretations and a date or scope. Thus this is an open-ended program, not a uniquely defined existence theorem. The supplied research summary explicitly treats it as such and points to related homotopy-type and virtual-knot work without claiming a resolution.
\subsection{Short English statement}
Seek additional geometric or topological ways to understand both the Jones polynomial and Khovanov homology. The source leaves the link category and standard for an interpretation open. The supplied assessment regards this as a continuing research direction rather than a single theorem with a defined endpoint.
\subsection{Sources}
\begin{itemize}
\item Local UnsolvedMath registry, permanent ID \(10600023\), exact source code \texttt{catalogue entry}, ``Virtual-knot problem 23.'' No individual proposer or proposed year is supplied.
\item \emph{Unsolved Problems in Virtual Knot Theory and Combinatorial Knot Theory}, revised survey, dedicated problem-list item 23: \url{https://arxiv.org/abs/1409.2823}. Background: 60 top-level items, compound items retained intact; source-TeX extraction, accessed 2026-07-29.
\item Supplied research context: the Lipshitz--Sarkar program on Khovanov homology and homotopy types, without a full citation; Rasmussen invariants for virtual knots (2016), arXiv:1603.xxxx; Khovanov homology of alternating virtual links (2019), arXiv:1904.xxxx. The xxxx identifiers are incomplete supplied placeholders and are not completed here.
\item Source evidence NEEDS\_REVIEW presents an open problem and requires a current primary-literature audit. Supplied OPEN-TRIAGE summary calls the request open-ended, reports no resolution expected or found, and says no single ``solution'' is available for this broad direction. This is a supplied scope assessment, not an independently established impossibility claim.
\item Dataset attribution: Ulam AI / UnsolvedMath Contributors, CC BY 4.0. Source rights NEEDS\_REVIEW: publicly accessible author/publisher copy, redistribution terms requiring review.
\end{itemize}
\end{document}
